\chapter[Równanie Schrödingera]{Równanie Schrödingera: hamiltoniany samosprzężone, ewolucja unitarna i dyspersja}
\label{chap:pde22}

\section{Od równania falowego do ewolucji unitarnej}
Niech $d$ będzie dodatnią liczbą całkowitą. W tym
rozdziale niewiadoma jest zespolona:
$u\colon\R\times\R^d\to\C$. Przyjmujemy jednostki, w których
$\hbar=1$ oraz masa cząstki jest równa $1/2$. Równanie
Schrödingera przyjmuje wówczas postać
\begin{equation}\label{eq:pde22-schrod}
 i\partial_tu=Hu,\qquad H=-\Delta+V(x),\qquad u(0)=u_0.
\end{equation}
Potencjał $V$ ma być rzeczywisty, chyba że wyraźnie podamy inne
założenia. Symbolem $H$ oznaczamy \emph{operator wraz z jego
dziedziną}, a nie tylko formalne wyrażenie różniczkowe.
Przez $\mathcal H$ oznaczamy przestrzeń Hilberta
$L^2(\R^d;\C)$. Przyjmujemy w niej iloczyn skalarny
$\langle f,g\rangle=\int_{\R^d}f(x)\overline{g(x)}\,\mathrm d\lambda_d(x)$,
liniowy w pierwszym argumencie.

Istotna różnica względem równania ciepła jest już widoczna w
transformacie Fouriera. Dla swobodnego równania $i u_t=-\Delta u$
mnożnik wynosi $e^{-it|\xi|^2}$: jego moduł jest stale równy $1$.
Nie ma więc tłumienia wysokich częstotliwości ani analogicznego
do ciepła wygładzania w skali Sobolewa. Za to oscylacja fazy
prowadzi do \emph{dyspersji} przestrzennej.

\begin{uwaga}[Porównanie trzech ewolucji]
\label{rem:pde-v43-evolution-contrast}
Równanie ciepła (rozdziały~\ref{chap:pde18} i~\ref{chap:pde20})
prowadzi dla nieujemnego operatora samosprzężonego
do mnożnika spektralnego $e^{-t\lambda}$, $t>0$,
który tłumi wysokie częstotliwości spektralne
(nie tylko wartości własne w przypadku widma ciągłego). Układ hiperboliczny z rozdziału~\ref{chap:pde21}
ma skończoną prędkość propagacji przy odpowiednich
założeniach. Swobodne równanie Schrödingera zachowuje
normę $L^2$, ale nie ma ogólnej skończonej prędkości
propagacji; samo zachowanie normy nie implikuje
zwartości ani wygładzania. Wszystkie trzy stwierdzenia
odnoszą się do \emph{różnych} operatorów i klas rozwiązań.
\end{uwaga}

\section{Samosprzężoność swobodnego hamiltonianu}
Używamy unitarnej normalizacji transformaty Fouriera z
rozdziału~\ref{chap:pde-fourier-Rd}:
\[
 \widehat f(\xi)=(2\pi)^{-d/2}
 \int_{\R^d}e^{-\mathrm i x\cdot\xi}f(x)\,\mathrm d\lambda_d(x).
\]
Wzór ten przyjmujemy początkowo dla funkcji
$f\in\mathcal S(\R^d)$ z przestrzeni Schwartza.
Twierdzenie Plancherela~\ref{thm:pde8-plancherel}
daje unitarne przedłużenie $\mathcal F$ na $\mathcal H$;
przez dualność przedłużamy je również na przestrzeń
dystrybucji temperowanych $\mathcal S'(\R^d)$.
Dla $s\in\R$ przyjmujemy przestrzeń Sobolewa
\[
 H^s(\R^d;\C)=
 \{f\in\mathcal S'(\R^d):
 (1+|\xi|^2)^{s/2}\widehat f\in L^2(\R^d;\C)\},
\]
z normą
$\|f\|_{H^s}^2=\int_{\R^d}(1+|\xi|^2)^s
|\widehat f(\xi)|^2\,\mathrm d\lambda_d(\xi)$.
Warunek oznacza, że ważona transformata jest
reprezentowana przez funkcję z $L^2$.
Dla całkowitego $s\ge0$ jest to przestrzeń słabo
różniczkowalnych funkcji z pochodnymi do rzędu $s$
w $L^2$, z równoważną normą: wynika to z
$\widehat{\partial^\alpha f}=(\mathrm i\xi)^\alpha\widehat f$
i porównania wielomianu
$(1+|\xi|^2)^s$ z sumą $\sum_{|\alpha|\le s}|\xi^\alpha|^2$.
Norma Fouriera będzie w tym rozdziale normą domyślną.

\begin{twierdzenie}[Swobodny hamiltonian]\label{thm:pde22-free-sa}
Niech $H_0\colon D(H_0)\subset L^2(\R^d;\C)\to L^2(\R^d;\C)$
będzie operatorem $H_0u=-\Delta u$ z dziedziną
$D(H_0)=H^2(\R^d;\C)$. Wtedy $H_0$ jest samosprzężony,
nieujemny i domknięty, a $\mathcal F H_0\mathcal F^{-1}$
jest operatorem mnożenia przez $|\xi|^2$ na dziedzinie
$\{g\in L^2:|\xi|^2g\in L^2\}$.
\end{twierdzenie}
\begin{proof}
Równość $\widehat{-\Delta u}=|\xi|^2\widehat u$
wynika z całkowania przez części dla $u\in\mathcal S$
i obowiązuje także dystrybucyjnie.
Dla $g\in L^2$ warunek $|\xi|^2g\in L^2$
jest równoważny warunkowi $(1+|\xi|^2)g\in L^2$:
implikacja w jedną stronę wynika z nierówności
trójkąta, a w drugą z $|\xi|^2\le1+|\xi|^2$.
Zatem $\mathcal F$ przenosi $H^2$ dokładnie na
dziedzinę operatora mnożenia
$Mg=mg$, $m(\xi)=|\xi|^2$.
Dziedzina ta jest gęsta w $L^2$, gdyż zawiera
wszystkie funkcje z $L^2$ o ograniczonym nośniku,
a $g\chi_{\{|\xi|\le n\}}\to g$ w $L^2$.
Operator jest symetryczny, ponieważ $m$ jest
rzeczywista i dla $g,h\in D(M)$ mamy
$\langle Mg,h\rangle=\langle g,Mh\rangle$.
W szczególności $M\subset M^*$.
Jeżeli $h\in D(M^*)$, to istnieje $k\in L^2$
takie, że $\langle Mg,h\rangle=\langle g,k\rangle$
dla wszystkich $g\in D(M)$. Dla każdego $R>0$
funkcja $g=\chi_{\{|\xi|\le R\}}(mh-k)$ należy
do $D(M)$, bo $m$ jest ograniczona na tej kuli.
Po podstawieniu otrzymujemy
\[
 \int_{\{|\xi|\le R\}}|mh-k|^2\,\mathrm d\lambda_d(\xi)=0.
\]
Biorąc kolejno $R\in\N$, dostajemy $mh=k$ prawie
wszędzie na $\R^d$. Zatem $h\in D(M)$ i $M^*h=Mh$,
co dowodzi przeciwnej inkluzji i samosprzężoności.
Nieujemność wynika z
\[
 \langle Mg,g\rangle
 =\int_{\R^d}m(\xi)|g(\xi)|^2\,\mathrm d\lambda_d(\xi)\ge0.
\]
Domkniętość można zobaczyć bezpośrednio: ze
zbieżności $g_n\to g$ i $mg_n\to h$ w $L^2$
wybieramy wspólny podciąg zbieżny prawie wszędzie
w obu przypadkach. W granicy otrzymujemy $h=mg$
prawie wszędzie, więc $g\in D(M)$ i $Mg=h$.
Unitarność $\mathcal F$ przenosi wszystkie te
własności na $H_0$.
\end{proof}

\begin{uwaga}[Dziedzina jest częścią twierdzenia]
Formalny laplasjan na $C_c^\infty(\Omega)$ dla obszaru
z brzegiem jest zwykle jedynie symetryczny, a nie samosprzężony.
Realizacje Dirichleta i Neumanna mogą dawać różne ewolucje.
Na $\R^d$ wybór $H^2$ jest kanoniczną samosprzężoną
realizacją swobodnego hamiltonianu. Jest ona
domknięciem operatora $-\Delta$ z dziedziną
$C_c^\infty(\R^d)$: jego norma grafowa
$(\|u\|_2^2+\|\Delta u\|_2^2)^{1/2}$ jest
równoważna normie $H^2$, ponieważ
$1+|\xi|^4\le(1+|\xi|^2)^2\le2(1+|\xi|^4)$,
a $C_c^\infty$ jest gęste w $H^2$.
W szczególności ta realizacja nie wymaga wyboru
dodatkowych warunków brzegowych.
\end{uwaga}

\section{Ewolucja unitarna i poprawność zagadnienia Cauchy'ego}
\begin{twierdzenie}[Grupa swobodnego Schrödingera]\label{thm:pde22-group}
Dla $t\in\R$ i $f\in L^2(\R^d;\C)$ niech
$U_0(t)f=\mathcal F^{-1}(e^{-it|\xi|^2}\widehat f)$.
Wtedy $(U_0(t))_{t\in\R}$ jest silnie ciągłą grupą
unitarną. Dla każdego $s\in\R$ ten sam wzór definiuje
silnie ciągłą grupę izometrii na $H^s$ (także
dla $s<0$, gdy $f$ nie musi należeć do $L^2$).
Dla każdego $f\in H^s$ zachodzi
$\|U_0(t)f\|_{H^s}=\|f\|_{H^s}$ oraz
$U_0(t)H^s=H^s$. Jej generator na $\mathcal H$
jest równy $-\mathrm iH_0$ z dziedziną $H^2$. Jeśli $f\in H^2$, to
$u(t)=U_0(t)f$ należy do
$C(\R;H^2)\cap C^1(\R;L^2)$ i spełnia
$i u_t=-\Delta u$. Jest to jedyne rozwiązanie w tej klasie
z daną $u(0)=f$.
\end{twierdzenie}
\begin{proof}
Ponieważ $|e^{-it|\xi|^2}|=1$, twierdzenie Plancherela
daje zachowanie iloczynu skalarnego w $L^2$,
a definicja normy Fouriera daje zachowanie normy
$H^s$ dla każdego $s\in\R$. Dla $s<0$ mnożnik
działa również na dystrybucje temperowane,
ponieważ jest gładki, a każda jego pochodna
ma co najwyżej wzrost wielomianowy.
Tożsamość mnożników
$e^{-i(t+r)|\xi|^2}=e^{-it|\xi|^2}e^{-ir|\xi|^2}$
daje prawo grupowe. Ponieważ $U_0(0)=I$,
operator $U_0(-t)$ jest odwrotnością $U_0(t)$,
więc izometrie są surjektywne.
Ponadto
$\|(U_0(t)-I)f\|_2^2=
\int_{\R^d}|e^{-it|\xi|^2}-1|^2|\widehat f(\xi)|^2\,\mathrm d\lambda_d(\xi)$.
Funkcja podcałkowa dąży punktowo do zera i jest ograniczona
przez $4|\widehat f|^2$; twierdzenie o zbieżności zmajoryzowanej
zapewnia silną ciągłość w zerze, a prawo grupowe w każdym $t$.
Tę samą argumentację można wykonać z wagą
$(1+|\xi|^2)^s$.
Jeśli $f\in H^2$, iloraz różnicowy w dziedzinie Fouriera
ma postać
$\frac{e^{-i(t+h)|\xi|^2}-e^{-it|\xi|^2}}h\widehat f$.
Dąży do $-i|\xi|^2e^{-it|\xi|^2}\widehat f$ i ma moduł
nie większy niż $|\xi|^2|\widehat f|$,
na mocy $|e^{\mathrm i y}-1|\le |y|$ dla $y\in\R$.
Kwadrat modułu różnicy ilorazu i jego granicy
ma całkowalną majorantę $4|\xi|^4|\widehat f|^2$.
Twierdzenie o zbieżności zmajoryzowanej daje
pochodną czasową $u_t=-\mathrm iH_0u$ w $L^2$.
Pochodna jest ciągła, gdyż $H_0:H^2\to L^2$
jest ograniczony i $u$ jest ciągła w $H^2$.
Odwrotnie, jeśli ilorazy
$(U_0(h)f-f)/h$ mają granicę w $L^2$ przy $h\to0$,
to ich transformaty wzdłuż pewnego ciągu $h_n\to0$
mają podciąg zbieżny prawie wszędzie do transformaty
tej granicy. Granica punktowa jest koniecznie
$-\mathrm i|\xi|^2\widehat f$, więc
$|\xi|^2\widehat f\in L^2$ i $f\in H^2$.
To identyfikuje dokładną dziedzinę generatora.
Dla dwóch rozwiązań
różnica $w$ spełnia to samo równanie z zerowymi danymi.
Różniczkując normę, uzyskujemy
$\frac{\mathrm d}{\mathrm dt}\|w\|_2^2=
2\operatorname{Re}\langle-iH_0w,w\rangle=0$,
ponieważ $\langle H_0w,w\rangle$ jest rzeczywiste.
Zatem $w=0$.
\end{proof}

\begin{stwierdzenie}[Granica regularności]
Dla każdej liczby $s\in\R$ i $t\in\R$ grupa $U_0(t)$
jest izometrią na $H^s(\R^d)$, ale dla $t\ne0$ nie
działa w ogólności jako operator wygładzający
$H^s\to H^{s+\eps}$ dla żadnego $\eps>0$.
\end{stwierdzenie}
\begin{proof}
Gdyby $U_0(t)f\in H^{s+\eps}$ dla każdego $f\in H^s$,
to ponieważ mnożnik odwrotny $U_0(-t)$ zachowuje
$H^{s+\eps}$, mielibyśmy $f=U_0(-t)U_0(t)f\in H^{s+\eps}$
dla wszystkich $f\in H^s$, co jest fałszem.
Dla ścisłości wybieramy dystrybucję temperowaną $f$
o transformacie danej funkcją
\[
 \widehat f(\xi)=
 \frac{\chi_{\{|\xi|\ge2\}}(\xi)}
 {(1+|\xi|^2)^{s/2+d/4}\log(2+|\xi|)}.
\]
Funkcja ta jest lokalnie całkowalna i ma co
najwyżej wzrost wielomianowy, więc określa
dystrybucję temperowaną. We współrzędnych
biegunowych kwadrat normy $H^s$ jest skończony,
gdyż całka radialna jest porównywalna z
$\int_2^\infty(r(\log r)^2)^{-1}\,\mathrm d\lambda_1(r)$,
która jest zbieżna przez podstawienie $y=\log r$.
Dla normy $H^{s+\eps}$ całka radialna jest
porównywalna z
$\int_2^\infty r^{2\eps-1}(\log r)^{-2}\,\mathrm d\lambda_1(r)$.
Jest rozbieżna, ponieważ dla dostatecznie dużego $r$
mamy $(\log r)^2\le r^\eps$, a całka
z $r^{\eps-1}$ jest rozbieżna.
Stąd $f\in H^s\setminus H^{s+\eps}$, a moduł
mnożnika $U_0(t)$ równy $1$ pokazuje również
bezpośrednio, że $U_0(t)f$ nie należy do
$H^{s+\eps}$.
\end{proof}

\begin{twierdzenie}[Zasada Duhamela]\label{thm:pde22-duhamel}
Niech $0<T<\infty$, $f\in L^1(0,T;\mathcal H)$ i
$u_0\in\mathcal H$. Funkcja
\[
 u(t)=U_0(t)u_0-i\int_0^tU_0(t-r)f(r)\,\mathrm d\lambda_1(r)
\]
należy do $C([0,T];\mathcal H)$ i jest jedynym
rozwiązaniem łagodnym $\mathrm i u_t=H_0u+f$
z warunkiem $u(0)=u_0$ (rozwiązanie łagodne
definiujemy jako funkcję ciągłą w $\mathcal H$
spełniającą powyższe równanie całkowe).
Spełnia również równanie w sensie dystrybucji
o wartościach w $H^{-2}$.
Spełnia oszacowanie
$\|u(t)\|_2\le\|u_0\|_2+\int_0^t\|f(r)\|_2\,\mathrm d\lambda_1(r)$.
Jeżeli ponadto $u_0\in H^2$ i
$f\in C([0,T];H^2)$, to rozwiązanie należy do
$C([0,T];H^2)\cap C^1([0,T];L^2)$ i spełnia równanie w $L^2$.
\end{twierdzenie}
\begin{proof}
Funkcja $r\mapsto U_0(-r)f(r)$ jest silnie
mierzalna w $\mathcal H$: dla funkcji prostych
wynika to z silnej ciągłości grupy, a dla ogólnej
$f$ przechodzimy do granicy prawie wszędzie
przez aproksymację funkcjami prostymi i izometryczność.
Jej norma jest równa $\|f(r)\|_2$,
więc jest całkowalna w sensie Bochnera.
Korzystając z prawa grupowego oraz ciągłości
operatora $U_0(t)$, zapisujemy rozwiązanie jako
\[
 u(t)=U_0(t)w(t),\qquad
 w(t)=u_0-\mathrm i\int_0^t U_0(-r)f(r)\,\mathrm d\lambda_1(r).
\]
Funkcja $w$ jest absolutnie ciągła w $\mathcal H$,
a $w'=-\mathrm i U_0(-t)f(t)$ prawie wszędzie.
Dla $t_n\to t$ mamy
\[
 \|U_0(t_n)w(t_n)-U_0(t)w(t)\|_2
 \le\|w(t_n)-w(t)\|_2+
       \|(U_0(t_n)-U_0(t))w(t)\|_2\to0.
\]
To daje ciągłość na całym $[0,T]$, warunek początkowy
i, przez nierówność dla normy całki Bochnera,
oszacowanie z tezy.

Operator $-\Delta$ ma ograniczone przedłużenie
$H_0:\mathcal H\to H^{-2}$, ponieważ
$|\xi|^2/(1+|\xi|^2)\le1$. Dla każdego $g\in\mathcal H$
ilorazy różnicowe $U_0(t)g$ zbiegają w $H^{-2}$
do $-\mathrm i H_0U_0(t)g$ (ten sam argument
zmajoryzowanej zbieżności z wagą $(1+|\xi|^2)^{-2}$).
Ponadto
$\|(U_0(t+h)-U_0(t))g\|_{H^{-2}}\le |h|\|g\|_2$.
Rozkładając różnicę wartości $u$ na składnik
z różnicą $w$ i składnik z różnicą $U_0$,
otrzymujemy absolutną ciągłość w $H^{-2}$:
pierwszy składnik jest kontrolowany przez
$\|w(t+h)-w(t)\|_2$, drugi przez $|h|\sup_t\|w(t)\|_2$.
W chwilach różniczkowalności $w$ w $\mathcal H$
reguła iloczynu, wynikająca z tego samego rozkładu,
daje w $H^{-2}$
\[
 u'=-\mathrm iH_0u+U_0(t)w'
    =-\mathrm i(H_0u+f)
\]
prawie wszędzie. Stąd równanie dystrybucyjne.

Przy mocniejszych danych $w$ jest klasy $C^1$
w $H^2$, dzięki silnej ciągłości grupy na $H^2$
i ciągłości $f$. Prawo grupowe daje ciągłość
$u$ w $H^2$. Rozkład ilorazu różnicowego teraz
w $\mathcal H$ i identyfikacja generatora dają
$u'=-\mathrm i(H_0u+f)$ w $\mathcal H$ we wszystkich
chwilach (na końcach jednostronnie).
Prawa strona jest ciągła w $\mathcal H$, co
dowodzi żądanej regularności.
Jednoznaczność łagodna wynika z definicji, a
jednoznaczność w klasie silnej z zachowania normy
dla różnicy dwóch rozwiązań.
\end{proof}

\section{Potencjały rzeczywiste i zachowanie energii}
\begin{twierdzenie}[Ograniczony potencjał rzeczywisty]\label{thm:pde22-potential}
Niech $V\in L^\infty(\R^d;\R)$ i niech $V$ działa jako
operator mnożenia na $L^2(\R^d;\C)$. Wtedy
$H=H_0+V$, $D(H)=H^2(\R^d)$, jest samosprzężony.
Grupa $U_V(t)=e^{-itH}$ jest unitarna i silnie ciągła.
Dla $u_0\in H^2$ rozwiązanie $u(t)=U_V(t)u_0$
zachowuje zarówno normę $L^2$, jak i energię
$E(u)=\langle Hu,u\rangle=
\int_{\R^d}(|\nabla u|^2+V|u|^2)\,\mathrm d\lambda_d$.
\end{twierdzenie}
\begin{proof}
Operator $V$ jest ograniczony i samosprzężony.
Suma $H_0+V$ jest symetryczna na $H^2$.
Pokażemy jej samosprzężoność bez przyjmowania
ogólnego twierdzenia o perturbacjach.
Dla $a>\|V\|_\infty$ mamy
$\|(H_0\pm ia)^{-1}\|_{L^2\to L^2}\le a^{-1}$,
co wprost wynika z mnożnika $(|\xi|^2\pm ia)^{-1}$.
Faktoryzacja
$H\pm ia=(I+V(H_0\pm ia)^{-1})(H_0\pm ia)$
pokazuje, że $H\pm ia$ są surjektywne, gdyż
$\|V(H_0\pm ia)^{-1}\|<1$ i odwrotność pierwszego
czynnika otrzymujemy z szeregu Neumanna.
Symetryczny operator o obu surjektywnych operatorach
$H\pm ia$ jest samosprzężony: istotnie dla
$z\in D(H^*)$ i $h=(H^*+ia)z$ dobieramy
$y\in D(H)$ takie, że $(H+ia)y=h$.
Wtedy $(H^*+ia)(z-y)=0$, a
$\ker(H^*+ia)=\operatorname{Ran}(H-ia)^\perp=\{0\}$;
zatem $z=y\in D(H)$. To dowodzi $H=H^*$.
Z twierdzenia spektralnego dla operatorów samosprzężonych
wynika istnienie silnie ciągłej grupy $e^{-itH}$ oraz
$U_V(t)D(H)=D(H)$, $HU_V(t)=U_V(t)H$.
Unitarność daje zachowanie normy, a ostatnia komutacja
prowadzi do
$\langle HU_V(t)u_0,U_V(t)u_0\rangle=
\langle Hu_0,u_0\rangle$. Całkowanie przez części
jest uzasadnione dla $u_0\in H^2$ i $V$ ograniczone.
\end{proof}

\begin{stwierdzenie}[Lokalny bilans prawdopodobieństwa]
Niech $V\in L^\infty(\R^d;\R)$ oraz
$u_0\in H^2(\R^d)$. Dla rozwiązania z
Twierdzenia~\ref{thm:pde22-potential} zachodzi w sensie
dystrybucji na $\R\times\R^d$
\[
 \partial_t|u|^2+\operatorname{div}j=0,
 \qquad j=2\operatorname{Im}(\overline u\,\nabla u).
\]
\end{stwierdzenie}
\begin{proof}
Ponieważ $u_t=-iHu=i\Delta u-iVu$,
$\partial_t|u|^2=2\operatorname{Re}(\overline u\,u_t)
=-2\operatorname{Im}(\overline u\,\Delta u)$.
Wyraz $-iV|u|^2$ nie wnosi wkładu do części rzeczywistej,
bo $V$ jest rzeczywisty. Z kolei
$\operatorname{Im}(\nabla\overline u\cdot\nabla u)=0$,
więc $\operatorname{Im}(\overline u\,\Delta u)
=\operatorname{div}\operatorname{Im}(\overline u\,\nabla u)$.
Tożsamość jest poprawna dla gładkich funkcji, a następnie
w sensie dystrybucji przez aproksymację w $H^2$.
\end{proof}


\begin{twierdzenie}[Zależny od czasu, ograniczony potencjał rzeczywisty]
\label{thm:pde22-time-potential}
Niech $T>0$, $V\in L^\infty((0,T)\times\R^d;\R)$,
$F\in L^1(0,T;L^2(\R^d;\C))$ oraz
$u_0\in L^2(\R^d;\C)$. Istnieje dokładnie jedna funkcja
$u\in C([0,T];L^2(\R^d;\C))$ spełniająca formułę Duhamela
\[
 u(t)=U_0(t)u_0-i\int_0^tU_0(t-s)
       \bigl(V(s,\cdot)u(s)+F(s)\bigr)\,\mathrm ds.
\]
Zachodzi oszacowanie
$\|u(t)\|_2\le\|u_0\|_2+\int_0^t\|F(s)\|_2\,\mathrm ds$.
Jeżeli $F=0$, to $\|u(t)\|_2=\|u_0\|_2$ dla każdego
$t\in[0,T]$, mimo że operator $H(t)=-\Delta+V(t,\cdot)$
może zależeć od czasu. Nie twierdzimy, że funkcja ta jest
rozwiązaniem silnym o wartościach w $H^2$.
\end{twierdzenie}
\begin{proof}
Stosujemy obraz oddziaływania $w(t)=U_0(-t)u(t)$.
Szukana równość jest równoważna równaniu całkowemu
\[
 w(t)=u_0-i\int_0^t\bigl(K(s)w(s)+G(s)\bigr)\,\mathrm ds,
 \quad K(s)=U_0(-s)M_{V(s)}U_0(s),\quad
 G(s)=U_0(-s)F(s).
\]
Operator $M_{V(s)}$ jest mnożeniem przez rzeczywistą
funkcję $V(s,\cdot)$. Dla prawie każdego $s$ operator $K(s)$
jest ograniczony, samosprzężony i
$\|K(s)\|\le\|V\|_\infty$. Funkcja
$s\mapsto K(s)z$ jest silnie mierzalna dla każdego $z\in L^2$;
wynika to z mierzalności $V$, separowalności $L^2$ i silnej
ciągłości $U_0(s)$.

Iteracje Picarda równania całkowego zbiegają jednostajnie
w $C([0,T];L^2)$: różnica kolejnych iteracji jest
ograniczona przez stałą pomnożoną przez
$(\|V\|_\infty T)^n/n!$. Dla dwóch rozwiązań nierówność
Gronwalla daje jednoznaczność. Ponadto
$w\in W^{1,1}(0,T;L^2)$ i prawie wszędzie
$w'=-iKw-iG$. Dla naszego iloczynu skalarnego, liniowego
w pierwszym argumencie, samosprzężoność $K$ daje
\[
 \frac{\mathrm d}{\mathrm dt}\|w(t)\|_2^2
 =2\operatorname{Re}\langle-iK(t)w(t)-iG(t),w(t)\rangle
 =2\operatorname{Re}\langle-iG(t),w(t)\rangle.
\]
Przy $G=0$ norma jest stała. Ogólnie, regularizując
pierwiastek $\sqrt{\|w\|_2^2+\delta^2}$ i przechodząc
$\delta\downarrow0$, otrzymujemy
$\|w(t)\|_2\le\|u_0\|_2+\int_0^t\|G(s)\|_2\,\mathrm ds$.
Unitarność $U_0(t)$ oraz $\|G(s)\|_2=\|F(s)\|_2$
kończy dowód.
\end{proof}


\begin{uwaga}[Potencjały nieograniczone]
Potencjały Coulomba i oscylatora harmonicznego wymagają
innych twierdzeń o samosprzężoności lub konstrukcji
przez formy kwadratowe. Nie wynikają one z powyższego
twierdzenia, ponieważ nie należą do $L^\infty$.
\end{uwaga}

\section{Jądro swobodnego propagatora i oszacowanie dyspersyjne}
\begin{twierdzenie}[Wzór na propagator i zanik dyspersyjny]
\label{thm:pde22-disp}
Dla $t\in\R\setminus\{0\}$ i $f\in\mathcal S(\R^d)$ zachodzi
\begin{equation}\label{eq:pde22-kernel}
 U_0(t)f(x)=(4\pi it)^{-d/2}
 \int_{\R^d}\exp\!\left(\frac{i|x-y|^2}{4t}\right)
 f(y)\,\mathrm d\lambda_d(y),
\end{equation}
gdzie potęgę zespoloną wybiera się zgodnie z ciągłą
granicą regularizacji gaussowskiej. W szczególności
\begin{equation}\label{eq:pde22-L1Linf}
 \|U_0(t)f\|_\infty\le(4\pi|t|)^{-d/2}\|f\|_1.
\end{equation}
Dla $2\le r\le\infty$ i $r'$ sprzężonego do $r$
zachodzi ponadto
\[
 \|U_0(t)f\|_{L^r}\le C_{d,r}
 |t|^{-d(1/2-1/r)}\|f\|_{L^{r'}}
\]
dla $f\in L^{r'}\cap L^2$, z przedłużeniem do wskazanych
klas przez gęstość.
\end{twierdzenie}
\begin{proof}
Wstawiamy czynnik $e^{-\eps|\xi|^2}$ do mnożnika
Fouriera i otrzymujemy, dla $\eps>0$, splot z funkcją
\[
 K_{t,\eps}(x)=[4\pi(\eps+it)]^{-d/2}
 \exp\!\left(-\frac{|x|^2}{4(\eps+it)}\right).
\]
Formuła wynika przez obliczenie transformaty Gaussa,
początkowo dla rzeczywistego parametru dodatniego,
a następnie przez analityczną kontynuację do półpłaszczyzny
$\operatorname{Re}(\eps+it)>0$. W dziedzinie Fouriera
mamy $e^{-(\eps+it)|\xi|^2}\widehat f\to
 e^{-it|\xi|^2}\widehat f$ w $L^2$ na mocy dominacji.
Dla $f\in\mathcal S$ i ustalonego $t\ne0$ dominacja
w całce przestrzennej daje także zbieżność punktową
$K_{t,\eps}*f\to K_t*f$, gdzie
$K_t(x)=(4\pi it)^{-d/2}\exp(i|x|^2/(4t))$.
Zbieżności muszą opisywać tę samą dystrybucję,
co dowodzi wzoru~\eqref{eq:pde22-kernel}.
Ponieważ $|K_t(x)|=(4\pi|t|)^{-d/2}$,
nierówność~\eqref{eq:pde22-L1Linf} jest natychmiastowa.
Drugie oszacowanie wynika przez interpolację
Riesza--Thorina między $L^1\to L^\infty$
a unitarnością $L^2\to L^2$. Przy parametrze
$\theta\in[0,1]$ mamy $1/r=\theta/2$
i $1/r'=1-\theta/2$; potęga czasu to
$-d(1-\theta)/2=-d(1/2-1/r)$.
\end{proof}

\begin{przyklad}[Pakiet gaussowski i brak stożka propagacji]
Dla $u_0(x)=e^{-|x|^2/2}$ otrzymujemy
\[
 U_0(t)u_0(x)=(1+2it)^{-d/2}
 \exp\!\left(-\frac{|x|^2}{2(1+2it)}\right).
\]
Dla każdego $t\ne0$ funkcja ta jest niezerowa we wszystkich
punktach $\R^d$. Zatem równanie Schrödingera, w przeciwieństwie
do równania falowego z rozdziału 21, nie ma skończonej prędkości
propagacji w sensie zachowywania zwartego nośnika.
Ściśle rzecz biorąc, pakiet Gaussa ma już w chwili początkowej
niezwarty nośnik. Poprawny kontrprzykład ze zwartymi danymi
otrzymujemy jednak ze wzoru~\eqref{eq:pde22-kernel}:
jeżeli $f\in C_c^\infty(\R^d)$ jest niezerowa, to dla $t\ne0$
$U_0(t)f$ jest, po wyłączeniu czynnika
$e^{i|x|^2/(4t)}$, transformatą Fouriera funkcji
o zwartym nośniku. Jest więc rzeczywiście analityczna
jako funkcja zmiennej rzeczywistej i nie może znikać na
niepustym zbiorze otwartym bez identycznego zerowania.
Unitarność wyklucza identyczne zerowanie. W szczególności
$U_0(t)f$ nie może mieć zwartego nośnika.
\end{przyklad}

\section[Oszacowania Strichartza w dwóch wymiarach]{Oszacowania Strichartza w modelowym przypadku $d=2$}
Oszacowanie dyspersyjne dotyczy ustalonego czasu. Pokażemy,
jak z niego uzyskać normę czasowo-przestrzenną.
Wybrany przypadek $(q,r)=(4,4)$ jest wystarczająco bogaty,
aby zobaczyć całą ideę metody $TT^*$, a zarazem pozwala
uniknąć najtrudniejszych zagadnień końcowych.

\begin{lemat}[Oszacowanie całkowe w czasie]\label{lem:pde22-hls}
Dla każdego $g\in L^{4/3}(\R)$ zachodzi
\begin{equation}\label{eq:pde22-hls}
 \left\|\int_{\R}|t-s|^{-1/2}|g(s)|\,\mathrm ds
 \right\|_{L^4_t}\le C\|g\|_{L^{4/3}_t}.
\end{equation}
\end{lemat}
\begin{proof}
Niech $Mg(t)=\sup_{R>0}(2R)^{-1}\int_{t-R}^{t+R}|g(s)|\,\mathrm ds$
będzie operatorem maksymalnym Hardy'ego--Littlewooda.
Najpierw przypomnimy wraz z uzasadnieniem jego ograniczoność
na $L^p(\R)$ dla $p>1$. Dla skończenie wielu przedziałów
wybranych z rodziny $\{I:|I|^{-1}\int_I|g|>\lambda\}$
można wybrać rozłączne $I_j$ tak, aby ich trzykrotne
powiększenia pokrywały całą rozważaną skończoną rodzinę.
Wynika stąd oszacowanie słabego typu
$|\{Mg>\lambda\}|\le C\|g\|_1/\lambda$
przez aproksymację otwartego zbioru poziomicowego.
Rozdzielmy teraz $g=g_1+g_2$,
$g_1=g\,\chi_{\{|g|>\lambda/2\}}$ oraz
$g_2=g\,\chi_{\{|g|\le\lambda/2\}}$.
Wtedy $Mg_2\le\lambda/2$, więc
$|\{Mg>\lambda\}|\le C\lambda^{-1}
\int_{\{|g|>\lambda/2\}}|g(s)|\,\mathrm ds$.
Po pomnożeniu przez $p\lambda^{p-1}$, całkowaniu
w $\lambda$ i zastosowaniu twierdzenia Tonellego
dostajemy $\|Mg\|_p^p\le C_p\|g\|_p^p$.

Dla ustalonego $R>0$ rozbijamy całkę w~\eqref{eq:pde22-hls}
na $|t-s|\le R$ i $|t-s|>R$. W pierwszej części
rozkład na pierścienie dyadyczne daje oszacowanie
$C R^{1/2}Mg(t)$. W drugiej nierówność Höldera z
wykładnikami $4$ i $4/3$ daje
$C R^{-1/4}\|g\|_{4/3}$, gdyż
$\int_{|r|>R}|r|^{-2}\,\mathrm dr=2/R$.
Wybieramy $R>0$ tak, aby
$R^{3/4}=\|g\|_{4/3}/Mg(t)$, jeśli oba człony
są dodatnie (przypadki zerowe otrzymujemy granicznie).
W rezultacie lewa strona całki punktowej jest ograniczona
przez $C(Mg(t))^{1/3}\|g\|_{4/3}^{2/3}$.
Podnosząc do czwartej potęgi i całkując, otrzymujemy
\[
 \|I_{1/2}|g|\|_4^4\le C\|Mg\|_{4/3}^{4/3}
 \|g\|_{4/3}^{8/3}\le C\|g\|_{4/3}^4.
\]
To dowodzi tezy.
\end{proof}

\begin{twierdzenie}[Oszacowanie Strichartza $L^4_{t,x}$ w $\R^2$]
\label{thm:pde22-strichartz}
Niech $d=2$, $u_0\in L^2(\R^2)$ i
$F\in L^{4/3}(\R\times\R^2)$. Wtedy
\begin{align}
 \|U_0(t)u_0\|_{L^4(\R\times\R^2)}&\le C\|u_0\|_2,
 \label{eq:pde22-str-hom}\\
 \left\|\int_{s<t}U_0(t-s)F(s)\,\mathrm ds
 \right\|_{L^4(\R\times\R^2)}&\le
 C\|F\|_{L^{4/3}(\R\times\R^2)}.
 \label{eq:pde22-str-inhom}
\end{align}
Drugą całkę definiujemy początkowo dla gładkich danych
o zwartym nośniku, a następnie przez domknięcie.
\end{twierdzenie}
\begin{proof}
Oznaczmy $T\varphi(t)=U_0(t)\varphi$ dla
$\varphi\in\mathcal S(\R^2)$. Formalnym operatorem
sprzężonym jest
$T^*G=\int_{\R}U_0(-s)G(s)\,\mathrm ds$
(dla $G\in C_c^\infty(\R\times\R^2)$).
Stąd
$TT^*G(t)=\int_{\R}U_0(t-s)G(s)\,\mathrm ds$.
Interpolacja oszacowania dyspersyjnego dla $d=2$
między $L^1\to L^\infty$ i $L^2\to L^2$ daje
$\|U_0(t-s)G(s)\|_{L^4_x}\le
 C|t-s|^{-1/2}\|G(s)\|_{L^{4/3}_x}$.
Z nierówności Minkowskiego oraz
Lematu~\ref{lem:pde22-hls} wynika zatem
\begin{equation}\label{eq:pde22-ttstar}
 \|TT^*G\|_{L^4_{t,x}}\le C\|G\|_{L^{4/3}_{t,x}}.
\end{equation}
Wykorzystujemy teraz dodatniość formy
$\langle TT^*G,G\rangle_{t,x}=\|T^*G\|_{L^2_x}^2$.
Nierówność Höldera wraz z~\eqref{eq:pde22-ttstar}
daje $\|T^*G\|_2^2\le C\|G\|_{4/3}^2$.
Dualność pomiędzy $L^{4/3}_{t,x}$ i $L^4_{t,x}$
implikuje zatem
$\|T\varphi\|_4\le C^{1/2}\|\varphi\|_2$.
Przez gęstość otrzymujemy~\eqref{eq:pde22-str-hom}.

Dla operatora opóźnionego zastępujemy całkę
$\int_{\R}$ przez $\int_{s<t}$. To samo oszacowanie
punktowe pozostaje prawdziwe, ponieważ
$\int_{s<t}|t-s|^{-1/2}g(s)\,\mathrm ds\le
\int_{\R}|t-s|^{-1/2}g(s)\,\mathrm ds$
dla nieujemnego $g$; Lemat~\ref{lem:pde22-hls}
prowadzi wprost do~\eqref{eq:pde22-str-inhom}.
\end{proof}

\begin{uwaga}[Admisybilność i zakres twierdzenia]
W wymiarze $d$ para wykładników $2\le q,r\le\infty$
jest formalnie zgodna ze skalowaniem, gdy
$2/q+d/r=d/2$. Powyżej \emph{nie} dowodziliśmy
wszystkich takich oszacowań ani przypadków końcowych.
W szczególności przypadek $(q,r)=(2,\infty)$ w wymiarze
$d=2$ nie spełnia standardowego twierdzenia Strichartza.
Oszacowania dla operatorów z potencjałem wymagają dodatkowych
założeń spektralnych; nie wynikają z samej samosprzężoności.
\end{uwaga}

\section{Stan związany i brak dyspersji dla pewnych hamiltonianów}
\begin{przyklad}[Oscylator harmoniczny]
W przestrzeni $L^2(\R^d)$ rozważmy nieujemny operator
$H_{\mathrm{osc}}=-\Delta+|x|^2$ w jego standardowej
samosprzężonej realizacji. Funkcja
$\psi_0(x)=e^{-|x|^2/2}$ spełnia
$H_{\mathrm{osc}}\psi_0=d\psi_0$.
W szczególności $e^{-itH_{\mathrm{osc}}}\psi_0
=e^{-itd}\psi_0$: jej gęstość przestrzenna
$|\psi_0(x)|^2$ nie zmienia się w czasie.
Nie może więc istnieć uniwersalne oszacowanie
$\|e^{-itH_{\mathrm{osc}}}f\|_\infty\le C|t|^{-d/2}\|f\|_1$
dla wszystkich $t\ne0$ i wszystkich $f\in L^1\cap L^2$,
bo zastosowanie go do $\psi_0$ i przejście $t\to\infty$
prowadzi do sprzeczności.
\end{przyklad}

\section{Zależność od warunków brzegowych i dalsze pytania}
Na ograniczonym obszarze gładkim laplasjan Dirichleta
ma dyskretne widmo. Jego ewolucja Schrödingera jest nadal
unitarna, lecz na funkcji własnej działa tylko przez zmianę
fazy, a zatem nie występuje globalny zanik dyspersyjny
analogiczny do~\eqref{eq:pde22-L1Linf}. Jest to kolejny
powód, dla którego należy wyraźnie oddzielać twierdzenia
na $\R^d$ od zagadnień z brzegiem.

Pozostają ważne kierunki dalszego rozwoju:
bezpośrednia konstrukcja hamiltonianów z osobliwymi potencjałami,
lokalne i globalne oszacowania Strichartza dla nieliniowych
równań Schrödingera, zjawiska rozpraszania oraz analiza
rezonansów spektralnych. Zagadnienia te wykraczają poza
podstawowy model liniowy niniejszego rozdziału.


